\documentclass[letterpaper,11pt]{article}
\usepackage[T1]{fontenc}
\usepackage[utf8]{inputenc}
\usepackage{lmodern}
\usepackage[french]{babel}
\usepackage[margin=1in,headheight=28pt]{geometry}
\renewcommand{\familydefault}{\sfdefault}
\usepackage{amsmath,amssymb,amsthm,bm,enumitem}
\usepackage[dvipsnames]{xcolor}
\usepackage[small,compact]{titlesec}
\usepackage{fancyhdr,fancyvrb}
\usepackage[colorlinks=true,allcolors=blue]{hyperref}
\pagestyle{fancy}
\fancyhead[L]{STT-2100 \\ Régression}
\fancyhead[R]{A.-A. Vallée \\ Université Laval}
\fancyfoot[C]{\thepage}
\setlength{\parindent}{0pt}
\setlength{\parskip}{5pt}
\fvset{fontsize=\small,frame=single,framesep=2mm}
\newcommand{\esp}{\mathbb{E}}
\newcommand{\var}{\operatorname{Var}}
\newcommand{\bo}[1]{\bm{#1}}
\newcommand{\N}{\mathcal{N}}
% Définir \avecsolutions avant l'input pour une version corrigée.
\newcommand{\prof}[1]{\ifdefined\avecsolutions\par{\color{blue}\textbf{Solution :} #1}\par\fi}
\begin{document}
\begin{center}{\Large Choix et validation de modèle avec R}\end{center}
Cette fiche accompagne les notes du chapitre 8. Les noms des objets peuvent être adaptés à vos données. Les exemples utilisent des données incluses dans R ou des données simulées; aucun fichier externe n’est requis.

\par\medskip\noindent\begin{minipage}{\linewidth}
\subsection*{Comparer des modèles candidats}
Les critères se comparent sur les mêmes observations et la même échelle de réponse. R compte aussi le paramètre de variance dans AIC() et BIC(). Une petite valeur est préférable.
\begin{Verbatim}
donnees <- na.omit(mtcars[c("mpg", "wt", "hp", "qsec")])
m1 <- lm(mpg ~ wt, data = donnees)
m2 <- lm(mpg ~ wt + hp, data = donnees)
m3 <- lm(mpg ~ wt + hp + qsec, data = donnees)
AIC(m1, m2, m3)
BIC(m1, m2, m3)
sapply(list(m1, m2, m3), function(m) summary(m)$adj.r.squared)
\end{Verbatim}

\end{minipage}\par
\par\medskip\noindent\begin{minipage}{\linewidth}
\subsection*{Vérifier les postulats}
Une courbure indique un problème de forme de la moyenne; un entonnoir suggère une variance non constante. Le graphique quantile-quantile examine la normalité. L'indépendance exige aussi de connaître le plan de collecte.
\begin{Verbatim}
par(mfrow = c(2, 2))
plot(m3)
par(mfrow = c(1, 1))
# PRESS : validation leave-one-out d'un modele fixe
press <- function(m) {
  h <- hatvalues(m)
  if (any(h >= 1 - 1e-10)) stop("Levier trop proche de 1")
  sum((residuals(m) / (1 - h))^2)
}
sapply(list(m1 = m1, m2 = m2, m3 = m3), press)
\end{Verbatim}

\end{minipage}\par
\par\medskip\noindent\begin{minipage}{\linewidth}
\subsection*{Validation croisée de modèles fixés à l’avance}
Utilisez les mêmes plis pour comparer les modèles. Ici les observations sont supposées indépendantes. Pour des données groupées ou temporelles, adaptez le découpage. Toute sélection ou transformation apprise doit être répétée dans l'entraînement de chaque pli.
\begin{Verbatim}
set.seed(2100)
K <- 5
pli <- sample(rep(seq_len(K), length.out = nrow(donnees)))
formules <- list(mpg ~ wt, mpg ~ wt + hp,
                 mpg ~ wt + hp + qsec)
eqm_cv <- function(formule) {
  erreurs <- numeric(nrow(donnees))
  for (k in seq_len(K)) {
    test <- pli == k
    ajuste <- lm(formule, data = donnees[!test, ])
    pred <- predict(ajuste, newdata = donnees[test, ])
    erreurs[test] <- (donnees$mpg[test] - pred)^2
  }
  mean(erreurs)
}
sapply(formules, eqm_cv)
\end{Verbatim}

\end{minipage}\par
\par\medskip\noindent\begin{minipage}{\linewidth}
\subsection*{Sélection algorithmique}
step() minimise un critère AIC; k=log(n) utilise une pénalité BIC. Cette recherche ne garantit pas le meilleur modèle parmi tous les sous-modèles. Les tests du modèle retenu ignorent la sélection; examinez stabilité, hiérarchie, diagnostics et performance externe.
\begin{Verbatim}
selection_aic <- step(m3, direction = "both", trace = 0)
selection_bic <- step(m3, direction = "both", trace = 0,
                      k = log(nrow(donnees)))
formula(selection_aic)
formula(selection_bic)
\end{Verbatim}

\end{minipage}\par
\par\medskip\noindent\begin{minipage}{\linewidth}
\subsection*{Aide}
Dans R, utilisez \texttt{?lm} et l’aide des fonctions présentées. 
\end{minipage}\par
\end{document}
